%============================================================
\chapter{Design of the Empirical Work: Connecting to Observation}
\label{ch:design}
%============================================================

\section{What is to be tested}

Part~\ref{part:theory} defined quantities and propositions, and Part~\ref{part:catalog}
deduced the subspaces that remain under constraints. Neither used observation. From this
part onwards they are connected to observation.

What is connected divides in two.

\begin{center}
\begin{tabularx}{\textwidth}{@{}lXl@{}}
\toprule
Object & Content & Question \\
\midrule
Quantities of Part~\ref{part:theory} & $\kappa$, $\CCC$, $g^\star$, $\phi$, $M(t)$, $\mathcal{G}$, and others & what values do they take \\
Deductions of Part~\ref{part:catalog} & the sieve of types, substitution, unmanned operation, collateral constraints & do they hold under the premises \\
\bottomrule
\end{tabularx}
\captionof{table}{The two objects the empirical work connects to.}
\label{tab:empirical-targets}
\end{center}

The first is measurement, the second verification. Because they need different preparation,
they are treated in stages below.

\section{Five stages to pass}
\label{sec:five-stages}

Connecting one quantity, or one proposition, to observation passes through five stages.
The organization is this text's own.

\begin{center}
\begin{tabularx}{\textwidth}{@{}llXl@{}}
\toprule
 & Stage & Question & What is needed if it fails \\
\midrule
1 & Operationalization & does the quantity have a unit and an object of observation & sharpening the definition \\
2 & Population & what set is to be enumerated & reselecting the object \\
3 & Measurement & can values be obtained at the granularity required & opening a route of acquisition \\
\midrule
4 & Reconciliation & does it agree with existing aggregates & --- \\
5 & Verification & are the implications fixed in advance supported & --- \\
\bottomrule
\end{tabularx}
\captionof{table}{Five stages from theory to observation.}
\label{tab:five-stages}
\end{center}

\textbf{The first three are prerequisites in series}: failing any one blocks the stages that
follow. \textbf{The last two are ways of using the values obtained, not a continuation of the
sequence.} They are separated by whether pre-registration has taken place.

The benefit of separating the stages is that the work required differs by where the process
stalls. A gap at stage 1 is not filled by collecting any amount of data, and a gap at stage 3
is not filled by refining any number of definitions.
\textbf{Without first judging where the process is stuck, effort is misallocated.}

\section{What to check at each stage}

\subsection{Stage 1: operationalization}

Three checks.

\begin{enumerate}[label=(\arabic*)]
\item Does the quantity have a \textbf{unit}?
\item Is the \textbf{object of observation} fixed --- the firm, a $\Phi$, or an identifier?
\item For a proposition, is the \textbf{sign uniquely fixed}?
\end{enumerate}

Some quantities of Part~\ref{part:theory} are operationalized by their definitions and some
are not. For the former, Chapter~\ref{ch:next} lists the variables to be recorded for each
$\Phi$.

For the latter --- quantities that are defined but whose unit or object of observation is
not fixed --- that has to be settled before proceeding to measurement.
\textbf{A defect at this stage cannot be detected by adding data.} Values do come out, so a
result appears without it being clear what was measured.

Item (3) is specific to propositions. If the theory does not fix the sign, either result is
consistent with it and no test is possible.

\subsection{Stage 2: population}

Three checks.

\begin{enumerate}[label=(\arabic*)]
\item Have the reference date, the scope and the lower bound been fixed?
\item What falls outside the sample, and is it correlated with the axis of interest?
\item Do objects that have stopped or exited remain in the sample?
\end{enumerate}

Items (2) and (3) are the central theme of this part and are treated in
Chapter~\ref{ch:method}. In this framework (3) weighs particularly heavily: the insolvency
time of Section~\ref{sec:objective} depends on $\Phi$, so an attempt to estimate $\Phi$
makes the composition of the sample itself a function of the object of estimation.

\subsection{Stage 3: measurement}

Three checks.

\begin{enumerate}[label=(\arabic*)]
\item Is the quantity measured anywhere?
\item If it is, can it be obtained at the granularity required?
\item If it cannot, which category of obstacle applies?
\end{enumerate}

The judgment at (3) governs the design of the survey, because the category determines
whether spending effort produces progress or nothing at all. Since the categories were
constructed through the survey work itself, they are placed in
Chapter~\ref{part:obstacles}.

\begin{remark}[Why this chapter does not supply the taxonomy itself]
\label{rem:taxonomy-placement}
The judgment at stage 3 needs the taxonomy of obstacles, but that taxonomy was constructed
through the empirical work. Placing it here would put findings not yet obtained ahead of
their basis.

\textbf{What this chapter supplies is the stages to check, not a classification of why a
stage was not passed.} The former is a design procedure, the latter a result of the survey.
The two can be stated independently: the stages say where the process stalled, the taxonomy
says why.
\end{remark}

\subsection{Stages 4 and 5: how the values obtained are used}

Passing stage 3 yields values. Their use divides in two, treated in the next section.

\section{Reconciliation versus verification}
\label{sec:comparison-vs-test}

Comparing a deduction with existing aggregates after the fact is called \emph{reconciliation}.
Judging an implication fixed in advance against data is called \emph{verification}.

\begin{center}
\begin{tabularx}{\textwidth}{@{}lXX@{}}
\toprule
 & Reconciliation & Verification \\
\midrule
When the implication is fixed & after & before \\
Declaring the population & not required & required \\
If they disagree & can be explained by another factor & the implication is rejected \\
Falsifiable & \textbf{no} & yes \\
\bottomrule
\end{tabularx}
\captionof{table}{Reconciliation versus verification.}
\label{tab:comparison-vs-test}
\end{center}

\textbf{Reconciliation cannot be falsified.} If the figures disagree one can say another
factor was at work, and nothing in the sample rules that explanation out. Reconciliation
therefore neither supports nor refutes the theory.

It still has a use. It confirms that a deduction is not wildly at odds with reality, and it
helps in choosing what to verify next.
\textbf{It has value as a check and none as evidence.}

The danger lies elsewhere.

\begin{remark}[Reconciliation can skip stages]
\label{rem:comparison-skips-stages}
Because reconciliation uses existing aggregates, \textbf{it works without declaring one's own
population}: stage 4 can be reached without passing stages 2 and 3.

This is the route by which reconciliation and verification are confused. An implication
supported after passing stages 1 to 3, and a direction that happens to agree with existing
aggregates after skipping stages 2 and 3, differ by an order of magnitude in the strength of
the evidence obtained. Outwardly both can be written up as ``consistent with the theory''.
\end{remark}

The comparison in Chapter~\ref{ch:solo} between the revival of family 3 and the fields of
business started by age is of this form. That section states explicitly that it is an
after-the-fact reconciliation and shows no causation.

Hence the following discipline.

\begin{quote}
\textbf{When a reconciliation is performed, say that it is one. Do not use it in place of an
implication fixed in advance.}
\end{quote}

This has the same form as the rule in Chapter~\ref{ch:method} that well-known firms are used
for illustration and not for inference. Both require that
\textbf{material not be used beyond the strength of the evidence it carries}.

\section{Where each stage is handled in this text}

\begin{center}
\begin{tabularx}{\textwidth}{@{}l>{\raggedright\arraybackslash}p{0.30\textwidth}X@{}}
\toprule
Stage & Where handled & Content \\
\midrule
1 Operationalization & Part~\ref{part:theory}, Chapter~\ref{ch:next} & the definitions, and the variables to record for each $\Phi$ \\
2 Population & Chapter~\ref{ch:method} & decomposition of selection bias and how to handle it \\
3 Measurement & Chapters~\ref{ch:instantiation} and~\ref{part:obstacles} & what could and could not be measured; the taxonomy of obstacles \\
4 Reconciliation & this chapter & the discipline of Section~\ref{sec:comparison-vs-test} \\
5 Verification & Chapter~\ref{ch:next}, Part~\ref{part:results} & the pre-registration procedure and the judgments on 22 implications \\
\bottomrule
\end{tabularx}
\captionof{table}{The five stages and where this text handles them.}
\label{tab:stage-mapping}
\end{center}

Chapter~\ref{ch:instantiation} opens Part~\ref{part:results} in order to tabulate the results
of stage 3. Making explicit which quantities could not be measured is that chapter's main
purpose, and that is a record of where between stages 1 and 3 the process stalled.

\begin{remark}[The stage reached differs by quantity]
\label{rem:stage-per-quantity}
The stage reached is not a single figure for the theory as a whole; it differs by quantity
and by proposition. Within one chapter, quantities that reached stage 5 sit alongside
quantities stalled at stage 1.

There is therefore no single answer to the question of how far the empirical work has come.
The stage reached can only be stated quantity by quantity.
\end{remark}

%============================================================

%---- Survey plan ----
\label{ch:next}

%============================================================

\section{Priorities for the survey}

The objects are narrowed to families 2 and 5. The selection can be derived deductively.

\paragraph{$\phi_{\mathrm{prod}}$ does not depend on $\Phi$.}

Of the three terms of~\eqref{eq:surplus}, $\phi_{\mathrm{prod}}=(p^{\ast}-c(\delta))\delta$
is fixed by marginal cost and the competitive price and contains nothing about the shape of
the map from delivery to settlement. Replacing $\Phi$ alone, holding $C_{\mathrm{prod}}$ and
marginal cost fixed, changes the sign of $\kappa$ and the working capital required, not the
production surplus (Table~\ref{tab:same-f-different-phi}).

\textbf{The shape of $\Phi$ bears on surplus through the other two terms.} To test the
framework, it suffices to choose $\Phi$ that isolate each of those two terms.

\paragraph{The two terms arise from different degrees of freedom.}

$\phi_{\mathrm{cog}}=p\cdot\E[\dbar-\delta]$ can be positive only when the contract states a
ceiling $\dbar$, that is, when $\pi$ does not refer to $\delta$
(Remark~\ref{prop:breakage}). That is degree of freedom (5). Relating the divergence to the
transaction frequency $\nu$ additionally requires the repetition of degree of freedom (6):
in a one-off, $\dbar-\delta$ is a single realization with no distribution, so no relation to
$\nu$ can be asked. \textbf{The family in which both (5) and (6) are present is family 2},
which Part~\ref{part:catalog} designates the principal home of cognitive surplus.

$\phi_{\mathrm{barg}}=(p-p^{\ast})\delta$ is a zero-sum transfer whose allocation is fixed by
whether the $y$ in $\pi=h(y)$ is verifiable. That is degree of freedom (4). In the order of
Table~\ref{tab:family-assignment-order},
\textbf{the only family decided by degree of freedom (4) is family 5}.

\paragraph{The two families sit diagonally in the quadrant diagram.}

Degrees of freedom (1) and (2) dominate, and those two axes give the four-way split of
Figure~\ref{fig:quadrant}. Family 2 is advance and flat ($\pi$ independent of $\delta$,
divergence possible), family 5 is arrears and linked ($\pi$ tied to $\delta$, no divergence):
\textbf{both axes take opposite values}. Choosing only one family would allow only one of
the two axes to be varied.

\paragraph{What can be tested.}

Family 2 permits direct tests of cognitive surplus (\eqref{eq:cog}) and of the capacity
proposition. In particular, that one and the same flat membership divides in two according to
whether a capacity constraint is present is the clearest instance of outwardly similar $\Phi$
having different durability.

Family 5 permits a direct test of the measurability condition (Chapter~\ref{sec:info}). The
claim that success fees and cost plus transfer risk in opposite directions is testable as a
correspondence between contract form and observability.

\paragraph{Why they are studied together.}

The two families rest on different institutional foundations, so data become available in
different ways. Family 2 has an exhaustive frame through registration under the Payment
Services Act; family 5 has per-firm filings published under the disclosure obligation of the
Employment Security Act.
\textbf{The former supplies a comprehensive frame, the latter individual variables.}

Looking at one family alone, when verification fails to progress
\textbf{one cannot tell whether to attribute it to a defect in the frame or a defect in the
variables}. Setting side by side two families with different foundations makes it possible to
separate which side is stuck. That is a further reason for taking both.

\section{Variables to record}

For each $\Phi$, at minimum the following are recorded.

\begin{center}
\begin{tabularx}{\textwidth}{@{}lX@{}}
\toprule
Variable & Content \\
\midrule
$\operatorname{sgn}\kappa$ & whether settlement precedes or follows delivery \\
$\CCC$ & only where a steady state holds; otherwise record that it does not \\
Form of dependence & $(F,p)$ of~\eqref{eq:two-part}, or $h(y)$ \\
Separation of $\dbar$ and $\delta$ & whether they separate, and an estimate of the divergence \\
Capacity constraint & whether present, and whether binding \\
$\nu$ & transaction frequency (\eqref{eq:decay}) \\
Measurability & whether effort or outcome is verifiable \\
Paying party & the beneficiary or a third party \\
Estimated split of $\phi$ & a qualitative allocation across $\phi_{\mathrm{prod}}$ / $\phi_{\mathrm{barg}}$ / $\phi_{\mathrm{cog}}$ \\
\bottomrule
\end{tabularx}
\captionof{table}{Minimum variables to record for each $\Phi$.}
\label{tab:variables-to-record}
\end{center}

\section{What must be pre-registered}

Following the conclusions of Part~\ref{part:method}, the following are fixed and recorded
before the survey begins.

\begin{enumerate}[label=(\arabic*)]
\item The definition of the population (reference date, geographic scope, industry scope, size floor)
\item The threshold for treating a state as steady (the criterion for stability of $r$)
\item How failure cases are collected, and the range in which they cannot be
\item Which of the APC effects is assumed zero, and on what grounds
\item The criterion distinguishing primary sources from retrospective accounts
\end{enumerate}

Deciding these after the fact lets the choices of Sections~\ref{sec:survivorship}
and~\ref{sec:disclosure} contaminate the conclusions.
